\chapter{Force from Geometry: the Inverse-Square Law}
\label{chap:force-from-geometry}

Part II began with a field that acts where it is not, and it ends with a law that no one needs to postulate
--- only to draw the surface that reveals it. The inverse-square fall-off --- the same $1/r^2$ that governs
gravity and the electric force --- is usually written down as a fact about force and left there. To the
instrument it is not a fact about force at all but a fact about geometry, and a simple one: a fixed budget of
influence, spread over a frontier that grows as a surface. The frontier is ours to draw --- where we set it,
how we shape it; the falloff it then forces, fixed by the dimension of the space it grows in, is not ours at
all. This single effect closes Part II, and despite its brevity it is one of the instrument's deep ideas,
because it shows two of the great forces of nature to be the same piece of bookkeeping wearing two names. With
it, the reading of a field as the bookkeeping a record carries between events is complete: the ledger ours to
keep, the balance the world's to enforce.

\section{The Inverse-Square Effect: influence over an expanding frontier}
\label{se:inverse-square}

Begin with what the instrument means by a field's reach. A source --- a mass, a charge --- has an influence it
sends outward, and the instrument records that influence as it crosses a \emph{frontier}: the set of distinct
directions, or slots, available at a radius $r$ from the source. In three dimensions that frontier is a
surface, and it grows quadratically with $r$. Discretely, the instrument counts a shell of about $6r^2$ slots
around the source; in the continuum the shell smooths into a sphere of area $4\pi r^2$. Either way, the
frontier at radius $2r$ is four times the frontier at $r$, and at $3r$ it is nine times: the surface a fixed
influence must cover swells as the square of the distance.

Now hold the budget fixed. The instrument's second premise is that the source's total influence is conserved
--- a charge does not gain or lose its reach as one walks away from it, and neither does a mass. Whatever the
source sends out, the same total crosses every frontier, near or far. But that fixed total is spread over a
frontier that keeps growing, so the share reaching any single slot must fall in exactly the proportion the
frontier rises,
\begin{equation}
  I(r) \;\propto\; \frac{1}{r^2}.
\end{equation}
There is the inverse-square law, and notice what produced it: not a force that happens to weaken with
distance, but a conserved budget divided by a growing surface. The budget is the world's to conserve, the
frontier ours to draw; the $1/r^2$ is only their quotient --- a constant over $r^2$, and the $r^2$ is the
frontier.

And the exponent is not arbitrary --- it is the dimension of space, read off the frontier. The surface a fixed
budget must cross grows as $r^2$ because space has three dimensions and the frontier at radius $r$ is a
two-dimensional sphere. Were space four-dimensional the frontier would be a three-sphere growing as $r^3$, and
the law an inverse cube; in $d$ dimensions the frontier grows as $r^{d-1}$ and the force falls as $1/r^{d-1}$.
The inverse \emph{square} is space announcing that it is three-dimensional: the $2$ in $1/r^2$ counts the
dimensions of the surface, one fewer than the dimensions of the space. The exponent is not ours to put
in by hand; the instrument reads it off how fast the frontier grows --- the dimension the world's, the surface
that shows it ours.

The simplest way to watch this is to count flux. Picture the source's influence as a fixed number of field
lines streaming outward, one bundle that never gains or loses a strand. Draw a sphere around the source and
count the lines crossing it; draw a larger sphere and count again. The count is the same --- every line that
leaves the source crosses every enclosing sphere exactly once, so the total flux through any surface around
the source is fixed at the budget. This is Gauss's law, and it is a statement about counts before it is an
equation:
\begin{equation}
  \oint \mathbf{E}\cdot d\mathbf{A} \;=\; \frac{Q}{\varepsilon_0},
\end{equation}
the same $Q$ on the right whatever sphere is on the left. The field --- the density of lines, the flux per
unit area --- is then the fixed count spread over the sphere's area,
\begin{equation}
  E \;=\; \frac{Q}{4\pi\varepsilon_0\, r^2},
\end{equation}
and this is read, not derived: the lines per area must fall as $1/r^2$ because the same lines cross an area
that grows as $r^2$. The surface is ours to draw; the count that crosses it is the world's. The instrument does
not solve Gauss's law for the field; it counts the lines, notes that the count is conserved, and reads the
dilution off the growing surface.

The field has an antiderivative, and it is gentler than the field. Sum the field's pull along a path inward
from far away and what accumulates is the potential, the work a unit of test budget owes to climb back out to
where it sits,
\begin{equation}
  \Phi(r) \;\propto\; \frac{1}{r},
\end{equation}
falling one power more slowly than the field that is its slope. The field is the potential's gradient,
$\mathbf{E} = -\nabla\Phi$, so the $1/r^2$ field is the downhill of a $1/r$ hill: the same conserved budget,
read once as a flux through the frontier and once as a height a test point has climbed --- two accounts of
one field, a flux and a height, and which book we keep is ours while the field they both record is the world's.
The potential is the smoother of the two, which is why the bookkeeping of energy so often runs through it
rather than the field.

The flux-counting has an exact form, the divergence theorem, and stating it makes the frontier picture
rigorous. For any field $\mathbf{F}$ and any region $\Omega$, the flux out through the boundary equals the
integral of the field's divergence over the interior,
\begin{equation}
  \oint_{\partial\Omega} \mathbf{F}\cdot d\mathbf{A} = \int_\Omega (\nabla\cdot\mathbf{F})\, dV .
\end{equation}
The divergence $\nabla\cdot\mathbf{F}$ is the density of sources --- how much the field springs from each point
--- and the theorem says the total springing-out, summed over the interior, equals the net flux escaping the
boundary. For the electric field this is Gauss's law in its local and global forms at once: the divergence
$\nabla\cdot\mathbf{E} = \rho/\varepsilon_0$ is the charge density, and its integral is the enclosed charge, so
the boundary flux is exactly the enclosed budget, $\oint \mathbf{E}\cdot d\mathbf{A} = Q/\varepsilon_0$,
however the boundary is drawn. The instrument's conserved budget over a growing frontier is this theorem read
backward: the flux is the same count $Q/\varepsilon_0$ through every enclosing surface because the interior
holds a fixed source, and the field thins as $1/r^2$ only because the frontier it crosses grows as $4\pi r^2$.

Here is the unification, and it is the chapter's real payoff. The instrument sees \emph{one} geometry behind
\emph{two} forces. Gravity pulls a test mass with $F = GMm/r^2$; the electric force pushes a test charge with
$F = k\,q_1 q_2/r^2$; and these are not two laws that happen to share an exponent. They are the same
frontier-dilution applied to two different budgets --- mass in the one case, charge in the other --- each
conserved, each spread over the same growing sphere. The instrument reading gravity and the instrument
reading electrostatics are doing the identical bookkeeping: a fixed budget over a frontier of area $4\pi r^2$.
The $1/r^2$ is shared not by coincidence but because the geometry is shared; change the budget and one
changes which force one is reading, but the law belongs to the surface, not to the force --- the budget ours to
name, the geometry the world's to enforce.

Two things still tell the budgets apart. Mass comes in one sign only, so its field lines have nowhere to end
but on other mass, and gravity only ever pulls; charge comes in two, so its lines run from positive to negative
and the force can push as readily as pull. The geometry is shared; the sign of the budget is not. And the
exponent itself has been weighed. Coulomb measured the inverse square for charge with a torsion balance,
Cavendish the same constant for gravity, and modern experiments pin the exponent to $2$ to many decimal places
--- a precision that matters, because any departure would betray a photon with mass, or a hidden extra
dimension changing how fast the frontier grows. The frontier is two-dimensional, the measurements say, to as
many places as they can reach, because space is three-dimensional to the same.

A second sighting confirms that the law lives in the geometry and not the source. Place a test point
\emph{inside} a uniform spherical shell of influence, and the net pull on it is exactly zero --- not
approximately, exactly --- wherever inside the shell it sits. The reason is the same symmetry that lit Arago's
spot: every patch of the shell pulling one way is balanced by the patch opposite it, the nearer patch smaller
and closer, the farther patch larger and more distant, and the two effects --- area growing as $r^2$,
influence falling as $1/r^2$ --- cancel term for term. The frontier's own geometry arranges its cancellation.
There is an even simpler way to see it, the same divergence theorem run on the inside: draw a Gaussian sphere
anywhere within the shell's hollow, and it encloses no source --- all the influence sits on the shell, outside
it --- so the flux through it is zero, $\oint \mathbf{E}\cdot d\mathbf{A} = 0$; and by the shell's spherical
symmetry the field can only point radially and depend on radius, so a vanishing flux forces the field itself to
vanish at every interior radius. Equivalently, the potential inside satisfies Laplace's equation with no
source, $\nabla^2\Phi = 0$, and a harmonic field with constant boundary values is constant throughout --- no
slope, no force. The interior is flat because the frontier put nothing there: the Gaussian sphere ours to draw anywhere in
the hollow, the exact zero it encloses the world's.
That a shell pulls nothing on its interior is the inverse-square law read from the inside, and it is one of
the plainest pieces of evidence that the $1/r^2$ is geometry and not force.

The honest floor is a smooth-shadow analogy. What the instrument actually counts is the discrete frontier ---
a shell of about $6r^2$ slots --- and the continuum sphere of area $4\pi r^2$ is the smooth shadow that
discrete shell casts as the slots refine. The bridge from ``a fixed budget shared over a growing frontier'' to
the continuum $1/r^2$ law is an analogy, named as one. The reading is that the inverse-square law is geometry,
not dynamics. Gravity and electrostatics dilute as $1/r^2$ for the same reason a fixed quantity of paint
spread over a growing sphere thins as $1/r^2$ --- the budget is conserved, the frontier expands, and the
quotient is the law.

This frontier returns, transformed, at the far end of the book. Here it is the surface a fixed influence is
diluted \emph{over}; in the relativity part it becomes the surface a region's information is bounded \emph{by}.
The holographic bound will say that the count a region can hold is set not by its volume but by the area of
its enclosing frontier --- the very $4\pi r^2$ that dilutes the force here, now capping the information there.
The inverse-square law and the holographic bound are the same frontier in two roles: once what a budget
thins across, once what a count cannot exceed. The instrument meets the frontier here, among the fields,
and will recognize it again at a black hole's horizon.

\bigskip

With the inverse-square law, Part II closes, and the field has been read off finite records from end to end. A
field, to the instrument, is what a record carries between events, and Part II has shown four of its faces. It
is a holonomy --- a residue a loop carries, surviving where the field itself is absent, and proved rather than
shadowed in the echo chamber. It is an interference --- a pattern that turns on whether the record can tell two
paths apart, the fringes the instrument's own ignorance made visible. It is a kernel --- the blind spot every
finite reading carries, the part of the world a projection is structurally unable to see. And it is a frontier
--- a budget diluted over a surface that grows as $r^2$, two forces revealed as one geometry. None of it
required a force imposed from outside; all of it was bookkeeping carried from one event to the next, read off
the count. Part III takes the same instrument to heat, entropy, and matter, where the bookkeeping becomes
statistical and the count, gathered over many events, becomes a temperature.

\bigskip
\begin{center}
\rule{0.4\textwidth}{0.4pt}
\end{center}
\bigskip

\noindent Testimony reaches only so far, and the way its weight thins with distance is not a failing of the
witness but a fact of geometry. A fixed measure of influence, sent out from a source, must cover a frontier
that widens as it goes --- and the same total, spread over an ever-larger surface, leaves less at any one place
the farther out that place sits. Nothing of the evidence is lost; it is only divided over more ground, thinning
as the square of the distance because the frontier grows as the square. Where we draw the frontier is ours; how
fast a fixed influence dilutes across it is the world's. The reach of a thing is the shape of the space it must
fill.
